\parthead{VIII}{A public research resource for Hong Kong}\label{part:8}
A historical dataset can be useful before it is complete, provided that incompleteness is visible. The immediate contribution of this project is an inspectable connection between source documents, dated geometry, a transparent mathematical model and reproducible calculations. It lets another reader challenge an assumption without having to guess where a number came from.

\fig{16_evidence_synthesis}{1}{The research's evidence structure. Historical observations and derived reconstructions occupy a different branch from assumed dynamics and constructed trajectories. Continuous-time model certification does not close the missing empirical connection. This is a dependency diagram, not a numerical measurement; no units apply. Sources: the project specification, settled obstacle readiness record and scenario result files.}{fig:synthesis}

The repository can support three concrete educational uses. Students can reproduce the bank--radius calculation and see why direction change takes time. A historical researcher can inspect different map editions and trace a proposed building identity back to its source. A controls researcher can replay a fixed witness, test a constraint or attempt a tighter bound without first accepting a claim of real-world superiority.

Responsible sharing keeps original sources, extracted observations and derivatives distinguishable. Lands Department open-data photographs and survey sheets are reproduced here under their attribution terms; Government intellectual-property ownership is retained. The data are supplied as-is without endorsement \citep{dataterms,csditerms}. Those permissions do not extend to all paid imagery or other institutions' documents. Accident-report scans and scholarly PDFs without established redistribution permission are linked and inventoried, not republished in this package.

The repository has not selected a blanket open-source licence for the original research and code. Public visibility is not permission to relicense third-party material. A broadly reusable resource will need a clearly chosen original-work licence alongside the existing source-specific terms.
\par\medskip
\subsection*{What a stronger public archive would contain}
The practical next step is a versioned evidence catalogue in which a map face, photograph, spot reading and flight-trace anchor each retain source identity, date precision, coordinate reference, extraction method and known limitations. Existing files already carry much of this information. A complete corridor-wide obstacle model and calibrated landing dataset remain future research deliverables, not hidden parts of this edition.

A reusable archive should keep three layers separable: original permitted data; derived measurements and reconstruction with reproducible transformations; and scenario calculations with explicitly frozen assumptions. Invalidated outputs must remain marked and traceable. A source update should not silently change the numerical meaning of an older published result.

Related research provides useful methodological comparisons. Wu and Chen study minimum-time Dubins-airplane paths in steady wind; Kim and coauthors formulate richer six-degree-of-freedom landing planning; FaSTrack explicitly separates planning from tracking-error guarantees \citep{wu,kim,fastrack}. These are pointers to established fields, not evidence that their assumptions or guarantees have been transferred to a historical 747. Their algorithms have not been independently reproduced for this white paper.

\subsection*{Returning to the original question}
The remarkable feature of Runway 13 was not a suspension of flight mechanics. A published instrument segment, visual references, a deliberate change of direction and skilled aircraft handling made use of a particular place. Mathematics helps explain the relation between bank, space, time, speed and stopping. Historical records explain why the place and the operating conditions cannot be reduced to a single elegant curve.

The constructed landings make one part of the answer concrete. Under declared reduced dynamics and constraints, the aircraft model can turn, descend, align, flare, touch down and stop. The references are explicit enough to inspect, and their inequalities are checked across complete time intervals. The historical reconstruction supplies dated evidence but has not yet closed the gap to a clearance-certified scene. CAL605 makes the ground phase and braking-system states impossible to ignore.

The next advance should be judged by whether it closes one of those identifiable gaps, not by whether a rendering looks more realistic or a numerical optimizer reports a larger score. A tighter family bound, a defensible vertical control, a recovered obstacle or a calibrated braking state would each make a distinct contribution.

That is the answer this edition can responsibly offer to the childhood question: the manoeuvre was possible because its demands could be coordinated, not because any one demand was absent. The remaining research question is how faithfully we can recover that coordination for a specified aircraft, date and set of conditions. The paper makes the current answer reproducible and the unanswered part precise.

\medskip{\small\textbf{Acknowledgement.} Aryan Ayyar initiated the project and directed its scientific scope. Research assistance, implementation, numerical checks, figures and drafting used AI tools. This edition does not imply independent human specialist review, institutional endorsement, or operational approval. The foreword remains explicitly for author review.}
\par\medskip
